\documentclass[11pt]{article}
\usepackage[margin=1in]{geometry}
\usepackage{graphicx}
\usepackage{amsmath,amssymb}
\usepackage{natbib}
\usepackage{booktabs}
\usepackage{xcolor}
\usepackage{hyperref}

\title{Strong Positive Selection Biases Identity-By-Descent-Based Inferences of Recent Demography and Population Structure in \textit{Plasmodium falciparum}}
\author{Anonymous}
\date{}

\begin{document}
\maketitle

\begin{abstract}
Malaria genomic surveillance often estimates parasite genetic relatedness using identity-by-descent (IBD), a summary that feeds effective population size ($N_e$), temporal demography, and network-based population structure inference. Yet strong recent positive selection, the default regime in \textit{Plasmodium falciparum} populations under antimalarial drug pressure, may bias IBD-based estimates in ways that have not been systematically evaluated. Using coalescent and forward-time simulations with a true-IBD inference algorithm, and whole-genome sequence data from 640 new and 4,026 publicly available \textit{P. falciparum} clinical isolates, we quantify this bias and evaluate a peak-removal correction. We show that a representative selection scenario (coefficient $s = 0.3$, selection start $80$ generations ago, single-origin favoured allele at $33.3$ cM on each chromosome) inflates long IBD segments and pairwise total IBD, underestimates recent $N_e$, and blurs InfoMap-detected community structure. Removing IBD segments that fall within validated peak regions partially restores both $N_e$ and structure in populations with low background relatedness, such as West Africa (WAF), but has little effect on estimates in populations with high background relatedness, such as Southeast Asia (SEA), where correction risks over-correction. We therefore recommend selection correction when IBD-based inferences are applied to high-transmission, low-relatedness settings and caution against routine correction in low-transmission, high-relatedness settings. Our simulation and true-IBD framework additionally provides a foundation for benchmarking IBD callers in high-recombination species.
\end{abstract}

\section{Introduction}

Malaria caused by \textit{Plasmodium falciparum} remains a leading cause of morbidity and mortality in tropical regions, and genomic surveillance has emerged as a central tool for monitoring transmission, tracking drug resistance, and guiding elimination efforts \citep{who2022worldmalaria,malariagen2021pf6}. A foundational primitive of this surveillance is genetic relatedness: the probability that two parasite genomes share recent common ancestry at a given locus. Identity-by-descent (IBD) segments, contiguous stretches of chromosome inherited from a recent common ancestor, summarise this relatedness at fine temporal resolution and have become a workhorse for inferring effective population size ($N_e$), recent demography, and fine-scale population structure in \textit{P. falciparum} and other recombining pathogens \citep{schaffner2018hmmibd,henden2018isorelate,taylor2019ibd}.

IBD-based inferences rest on an underlying coalescent or Wright--Fisher model in which the distribution of segment lengths maps directly to the time to most recent common ancestor (TMRCA) and in which pairwise total IBD reflects genome-wide shared ancestry. The flagship $N_e$ estimator in this family, \texttt{IBDNe}, uses the empirical distribution of segment lengths to reconstruct $N_e$ trajectories over the past dozens of generations and has been widely applied to human and parasite populations \citep{browning2015ibdne,daniels2015modeling}. Parallel lines of work have repurposed IBD for community detection on relatedness networks \citep{rosvall2008maps}, for selection scans via positional enrichment statistics such as $X_{iR,s}$ \citep{henden2018xirs}, and for connecting population genetic structure to epidemiologic transmission \citep{chang2019mapping,wong2017genetic}.

The critical assumption shared by these methods is that the fraction of the genome shaped by natural selection is small relative to the neutral background, so that the IBD distribution primarily reflects demography. In \textit{P. falciparum}, however, this assumption is routinely violated. Parasites in Southeast Asia and, increasingly, in Africa, carry multiple simultaneous selective sweeps driven by antimalarial drugs: the \textit{kelch13} locus for artemisinin resistance \citep{ariey2014kelch13,ashley2014tracking}, \textit{pfcrt} for chloroquine and piperaquine \citep{fidock2000pfcrt,amato2017piperaquine}, and a genetic background of accessory loci including \textit{arps10}, \textit{pph}, ferredoxin, and \textit{pfmdr1} \citep{miotto2015genetic}. Selective sweeps elongate the IBD segments that traverse the selected locus, inflate pairwise total IBD, and concentrate a disproportionate share of shared ancestry in a few chromosomal regions. Whether and how this distortion propagates into $N_e$ estimates and network-based structure inference has not been characterised, and there is no established correction protocol for IBD-based analyses in strongly selected populations.

A related literature on LD-based $N_e$ estimation has arrived at mixed conclusions. Classical theory predicts that selection shifts the LD decay curve and therefore biases $N_e$, yet \citet{novo2022ld} recently reported that LD-based estimators are virtually robust to natural selection across a broad range of parameters. Their simulations, however, used weak selection coefficients, randomly distributed selected loci, and long time scales, which capture a human-like regime rather than the concentrated, recent, strong selection typical of \textit{P. falciparum} under drug pressure. This apparent robustness therefore does not imply that IBD-based estimators, which depend on segment-length tails in a stronger way, inherit the same property.

In this work we directly quantify the effect of strong recent positive selection on IBD-based inferences in \textit{P. falciparum} and evaluate a simple correction based on IBD peak removal. Our contributions are fourfold. First, we assemble whole-genome sequencing data from 640 new and 4,026 publicly available clinical isolates, yielding an analyzable Southeast Asian (SEA) dataset of 2,055 parasites distributed across 14 years and 18 provinces in Cambodia, Thailand, Laos, and Vietnam, and a validation West African (WAF) dataset from the \textit{P. falciparum} catalogue \citep{malariagen2021pf6,manske2012analysis}. Second, we introduce \texttt{tskibd}, a true-IBD algorithm that operates on simulated genealogical trees in tree-sequence format \citep{kelleher2018efficient,haller2019treeseq,baumdicker2022msprime} and therefore decouples questions about IBD quality from questions about selection bias. Third, we use single-population and multi-deme stepping-stone simulations \citep{haller2019slim,kelleher2016msprime} to show that a representative selection scenario ($s = 0.3$, selection beginning $80$ generations ago, single-origin favoured allele at $33.3$ cM on each chromosome) distorts segment-length, pairwise total, and chromosomal positional IBD distributions, underestimates recent $N_e$ via \texttt{IBDNe}, and blurs InfoMap-detected community structure. Fourth, we develop and evaluate an IBD peak removal strategy, showing that the correction recovers accurate $N_e$ and structure in low-background-relatedness populations (WAF) but has little effect and risks over-correction in high-background-relatedness populations (SEA) \citep{daniels2015modeling,verity2020diversity}.

\section{Related Work}

\paragraph{IBD inference in \textit{P. falciparum}.} Two hidden Markov model callers dominate haploid parasite IBD inference: \texttt{hmmIBD} \citep{schaffner2018hmmibd} and \texttt{isoRelate} \citep{henden2018isorelate}. Both emit a binary per-site IBD state along pairs of haplotypes and summarise relatedness via the fraction of sites called IBD. Segment-level properties, although reported, are not routinely validated. \citet{taylor2019ibd} present a Bayesian framework that integrates relatedness uncertainty for downstream estimators.

\paragraph{Demographic estimators.} \texttt{IBDNe} \citep{browning2015ibdne} uses the empirical distribution of IBD segment lengths to estimate $N_e$ over the recent past, and has been deployed on \textit{P. falciparum} data from high- and low-transmission settings \citep{daniels2015modeling,wong2017genetic}. LD-based alternatives place different weight on recombination-distance bins and have been argued to be robust to selection in simulated human-like settings \citep{novo2022ld}; this claim is the closest analogue to ours but reaches a different conclusion under a different simulation regime, which we examine in the Discussion.

\paragraph{Selection detection with IBD.} IBD-based selection scans exploit positional enrichment of segments near favoured alleles. The $X_{iR,s}$ statistic of \citet{henden2018xirs} summarises this enrichment relative to a null expectation and motivates our peak-identification procedure, which applies a chromosome-wise threshold and a boundary-extension heuristic to flag sweep regions for removal.

\paragraph{Tree-sequence simulation.} Our true-IBD pipeline builds on modern tree-sequence tooling: \texttt{msprime} \citep{kelleher2016msprime,baumdicker2022msprime} for coalescent backgrounds, \texttt{SLiM} \citep{haller2019slim} with tree-sequence recording \citep{haller2019treeseq,kelleher2018efficient} for selection and structured demographies, and \texttt{tskit} data structures throughout. Recording the genealogy in simulation lets us bypass the confounding between IBD caller accuracy and the selection effects we wish to isolate.

\paragraph{Network community detection.} We use InfoMap \citep{rosvall2008maps} for community detection on IBD-weighted networks, evaluated against ground-truth labels via modularity \citep{newman2004modularity} and the adjusted Rand index. Within-community and inter-community IBD ratios characterise the strength of population boundaries before and after peak correction.

\paragraph{Empirical context.} Genomic surveillance in \textit{P. falciparum} has documented the spread of artemisinin resistance across the Greater Mekong Subregion \citep{ashley2014tracking,miotto2015genetic} and the sustained high relatedness of SEA parasite populations as transmission declined \citep{amato2017piperaquine}. West Africa, by contrast, has higher $N_e$ and lower background relatedness, even as Kelch13 variants have emerged locally \citep{verity2020diversity}. The contrast between these settings is central to our finding that background relatedness is the key modifier of selection-induced IBD bias and the utility of correction.

\section{Methods}

\subsection{Datasets}

We assembled a Southeast Asian (SEA) dataset of 2,055 analyzable \textit{P. falciparum} clinical isolates after quality control, combining 751 in-house isolates (640 newly sequenced here) from Cambodia and Thailand with 1,304 eastern SEA isolates from the MalariaGEN Catalogue of Genetic Variation in \textit{P. falciparum} v6.0 (Pf6) \citep{malariagen2021pf6}. Isolates span $14$ years and $18$ provinces across Cambodia, Thailand, Laos, and Vietnam. We obtained a West African (WAF) validation set from the same Pf6 release to represent a high-transmission, low-relatedness setting. Per-isolate coverage was evaluated as the fraction of the Pf genome sequenced at $\geq 5\times$, $\geq 10\times$, and $\geq 25\times$, and within-host complexity was estimated using the $F_{ws}$ statistic, with monoclonality defined as $F_{ws} > 0.95$. For polyclonal isolates, the predominant haploid genome was inferred using \texttt{dEploid} \citep{zhu2018deploid} when a dominant clone was detectable (ratio $<1$); otherwise, the isolate was excluded.

\subsection{True-IBD inference in simulation: tskibd}

Because candidate IBD callers have unknown error profiles under strong selection, we implemented \texttt{tskibd}, a true-IBD inference algorithm that reads simulated genealogical trees in tree-sequence format. For a pair of haploid samples, \texttt{tskibd} scans the genome, identifies every contiguous interval along which the two samples descend from a common ancestor younger than a configurable TMRCA cutoff, and emits these intervals as true IBD segments. This bypasses phased-genotype HMMs and provides a ground-truth oracle against which downstream summaries can be evaluated. We validated \texttt{tskibd} by reconstructing $N_e$ in neutral simulations under multiple demographic scenarios and comparing to the known population size.

\subsection{Simulation models}

We used two classes of simulations. The single-population model simulates a panmictic population under neutrality or with positive selection at one specified locus per chromosome; this model is used to isolate the effect of selection on the IBD segment-length distribution, on pairwise total IBD, and on chromosomal positional IBD. The multi-population model is a one-dimensional stepping-stone with five subpopulations (denoted $p_1, \ldots, p_5$) split from an ancestral population, with symmetric migration between adjacent demes at rate $m = 0.01$ per generation; a favoured allele is introduced into the leftmost deme and allowed to spread along the chain. Forward-time selection was implemented using \texttt{SLiM} with tree-sequence recording \citep{haller2019slim,haller2019treeseq}, and neutral coalescent backgrounds were generated with \texttt{msprime} \citep{kelleher2016msprime,baumdicker2022msprime}. Unless stated otherwise, representative results use selection coefficient $s = 0.3$, selection starting $80$ generations ago, and a single origin of the favoured allele introduced at $33.3$ cM of each chromosome.

\subsection{IBD-based estimators}

For empirical data, we called IBD with \texttt{hmmIBD} \citep{schaffner2018hmmibd}. $N_e$ trajectories were inferred with \texttt{IBDNe} \citep{browning2015ibdne}. Short IBD segments (0.2--2 cM) were included to cover the more distant past ($>25$ generations ago). IBD-based population structure was inferred by constructing weighted networks whose nodes are isolates and whose edge weights are pairwise genome-wide total IBD; we applied InfoMap community detection \citep{rosvall2008maps} and evaluated modularity against true or geographic labels \citep{newman2004modularity}. For community comparison, we used the adjusted Rand index.

\subsection{IBD peak identification and removal}

We identified IBD peak regions per chromosome using a threshold-based method on the IBD coverage profile, and validated candidate peaks using the $X_{iR,s}$ selection statistic \citep{henden2018xirs}. To capture neutral loci linked to the selected site, peak boundaries were extended outward to the chromosome-wise median line. Any IBD segment falling within a validated peak region was removed before re-running $N_e$ and structure estimators.

\subsection{Statistical tests}

Signed-rank tests were applied across replicated simulations to assess the significance of selection effects on $N_e$, with Bonferroni adjustment for multiple comparisons; differences at adjusted $p < 0.05$ were considered significant. For empirical community comparisons between pre- and post-peak-removal inferences, jackknife resampling was used to derive confidence intervals on the adjusted Rand index.

\section{Results}

\subsection{Parasite isolates and whole-genome sequencing data summary}

The SEA cohort of 2,055 analyzable isolates includes 751 (640 new) in-house isolates from Cambodia and Thailand and 1,304 eastern SEA isolates from Pf6, distributed across $14$ years and $18$ provinces in Cambodia, Thailand, Laos, and Vietnam. Among these isolates, $79.3\%$, $68.0\%$, and $46.1\%$ had at least $5\times$, $10\times$, and $25\times$ coverage across $>80\%$ of the \textit{P. falciparum} genome, respectively. The $F_{ws}$ statistic classified $80\%$ of isolates as monoclonal ($F_{ws} > 0.95$); among the remaining polyclonal isolates, $44.3\%$ harboured a predominant clone, which was included in the analyzable dataset. The WAF validation set exhibited a lower monoclonal fraction of $50.7\%$, consistent with the higher multiplicity of infection expected in a high-transmission setting. Figure~\ref{fig:dataset} summarises these cohort statistics.

\begin{figure}[t]
  \centering
  \includegraphics[width=0.95\linewidth]{figures/fig1_dataset.pdf}
  \caption{Summary of \textit{P. falciparum} isolates and whole-genome sequencing data used in this study. Left: distribution of per-isolate coverage among SEA isolates, with $79.3\%$, $68.0\%$, and $46.1\%$ covered at $\geq 5\times$, $\geq 10\times$, and $\geq 25\times$ over more than $80\%$ of the genome, respectively. Right: clonality summary; $80\%$ of SEA isolates are monoclonal by $F_{ws} > 0.95$, $44.3\%$ of polyclonal SEA isolates carry a predominant clone, and $50.7\%$ of WAF isolates are monoclonal.}
  \label{fig:dataset}
\end{figure}

\subsection{Positive selection distorts the IBD distribution and underestimates $N_e$}

Strong recent positive selection distorts all three major summaries of the IBD distribution and biases $N_e$ downward in recent generations. To isolate this effect cleanly, we used our single-population simulation paired with \texttt{tskibd} to derive true IBD under neutrality and under a representative selection scenario ($s = 0.3$, selection starting $80$ generations ago, single-origin favoured allele at $33.3$ cM on each chromosome); this setup removes confounders from IBD caller error because the genealogy is known exactly. Three aspects of the IBD distribution changed systematically under selection. First, the proportion of long IBD segments increased, pulling the segment-length distribution rightward. Second, pairwise genome-wide total IBD sharing between isolate pairs increased, yielding more closely related pairs in aggregate. Third, positional IBD coverage peaked near selected sites, yielding easily identifiable enrichment in chromosomal sweep regions. Because the \texttt{IBDNe} estimator derives recent $N_e$ from the segment-length distribution, and because selection disproportionately populates the long tail, the inferred recent $N_e$ is lower than the true population size.

We then tested whether removing IBD segments located within validated peak regions can correct this bias. Applying the peak identification and removal procedure to the selection simulation largely restored the segment-length distribution and yielded $N_e$ trajectories that closely match both the neutral simulation and the known true population size. Short segments in the range $0.2$--$2$ cM, which carry information about the more distant past ($>25$ generations ago), were retained throughout. Signed-rank tests across replicated simulations confirmed that the effect of selection on $N_e$ estimates is statistically significant at Bonferroni-adjusted $p < 0.05$. Across varying selection coefficients, intermediate selection durations, and small numbers of origins (hard sweeps), selection was stronger when the favoured allele had a greater chance of establishment, and the IBD-based $N_e$ bias tracked establishment probability. Figure~\ref{fig:ne_schematic} summarises these trajectories.

\begin{figure}[t]
  \centering
  \includegraphics[width=0.85\linewidth]{figures/fig2_ne_schematic.pdf}
  \caption{Schematic effect of strong recent positive selection on $N_e$ inference. Trajectories are anchored to the reported SEA range of $N_e \sim 10^3$ to $\sim 10^4$ over the most recent $60$--$80$ generations and to the reported selection parameters ($s = 0.3$, selection start $80$ generations ago, single-origin favoured allele at $33.3$ cM). Selection depresses inferred recent $N_e$, and IBD peak removal restores the trajectory toward the neutral/true curve.}
  \label{fig:ne_schematic}
\end{figure}

\subsection{Selection blurs IBD-based population structure}

Selection not only distorts $N_e$ but also degrades the signal of spatial population structure in IBD networks. To quantify this, we used the multi-deme stepping-stone model (Figure~\ref{fig:stepping_stone}) with five subpopulations and migration rate $m = 0.01$ between adjacent demes, which mimics the allele-frequency gradient expected when a selective sweep propagates through a structured parasite population. Under neutrality, within-population IBD sharing dominates the pairwise heatmap, the modularity of the IBD network with respect to true subpopulation labels is high, and InfoMap community detection recovers the true labels with high consistency. Under strong selection ($s = 0.3$), both within- and between-population IBD sharing increase, but the ratio of inter-population to intra-population IBD grows. Network modularity drops and InfoMap communities collapse, making adjacent demes difficult to distinguish. After removing IBD segments within validated peaks, the dominance of within-population sharing is restored, modularity rebounds, and InfoMap again recovers the true labels. This pattern was robust across a range of selection coefficients and replicated simulations, and the blurring effect is selection-strength-dependent.

\begin{figure}[t]
  \centering
  \includegraphics[width=0.90\linewidth]{figures/fig3_stepping_stone.pdf}
  \caption{One-dimensional stepping-stone model used to study selection effects on population-structure inference. Five subpopulations $p_1, \ldots, p_5$ are connected by symmetric migration between adjacent demes at rate $m = 0.01$. A favoured allele is introduced at the left deme and propagates along the chain, generating an allele-frequency gradient across subpopulations.}
  \label{fig:stepping_stone}
\end{figure}

\subsection{Genome-wide IBD sharing and selection signals in SEA isolates}

Empirical IBD coverage in SEA peaks precisely at known drug-resistance loci, confirming that the selection signature our simulations predict is present in the real data. We therefore asked whether selection signatures observed in simulation are visible in empirical SEA isolates. Using \texttt{hmmIBD} on the 2,055 SEA isolates, we plotted chromosomal IBD coverage and overlaid the positions of known or putative drug-resistance genes. IBD coverage peaks coincided with \textit{pfmdr1}, \textit{pfaat1}, \textit{pfcrt}, \textit{dhps}, \textit{pph} (PF3D7\_1012700), \textit{gch1}, \textit{kelch13}, and \textit{arps10}, and with \textit{ap2-g}/\textit{ap2-g2}, which have been implicated in altered sexual investment and increased transmission of resistant parasites \citep{ariey2014kelch13,fidock2000pfcrt,miotto2015genetic}. Including all parasite genomes, we observed a high baseline IBD across the genome compared to WAF, consistent with the decline in malaria transmission in SEA following intensive elimination efforts. To avoid confounding by high background relatedness, we iteratively removed highly related isolates, defined as those with IBD sharing greater than half the genome to many others, yielding an unrelated subset of $n = 701$ isolates whose baseline IBD is roughly five times lower. Peaks at \textit{pfcrt} and \textit{arps10} were substantially easier to identify in the unrelated subset, motivating the use of this subset for downstream $N_e$ and structure inference.

\subsection{IBD-based inference in SEA: low transmission, high background relatedness}

In SEA, IBD-based $N_e$ declined sharply over recent generations and population structure was dominated by geographically and resistance-stratified communities, but peak-region removal changed none of these inferences materially. On the $n = 701$ unrelated SEA isolates, \texttt{IBDNe} inferred a decreasing $N_e$ trajectory from approximately $10^4$ to approximately $10^3$ over the most recent $60$--$80$ generations, consistent with the rapid decline in malaria incidence in the region. InfoMap community detection identified five communities with more than $20$ isolates. The largest community, C0, was enriched for parasites sampled from Western Cambodia (Battambang, Kratie, Oddar Meanchey, and Pursat). C1, the second largest, spanned Northeastern Cambodia (Ratanakiri), Laos, and Vietnam, and had a lower within-community IBD level, reflecting distant relatedness. Hierarchical clustering of the community-level average IBD matrix placed C2, C3, and C4 closer to C0 than to C1. Temporal dynamics included the convergence of Pursat parasites into a main community in recent years and a similar shift in Oddar Meanchey $2$--$3$ years later, consistent with the documented spread of artemisinin-resistant lineages from west to north.

The drug-resistance mutation landscape of these communities was heterogeneous. The C0/C2/C3/C4 group exhibited elevated frequencies of resistance-associated or resistance-linked alleles, including PF3D7\_0720700 C1484F, ARPS10 V127M, PfCRT I356T, Ferredoxin D193Y, Kelch13 C493H, R539T, and C580Y, PPH V1157L, DHFR I164L, DHPS A581G (antifolate resistance), and PfCRT H97Y, I218F, and G353V (piperaquine resistance). C1, by contrast, showed low or no frequency at these loci and lacked Kelch13 mutations, consistent with the presence of multiple artemisinin-resistant founder populations and an artemisinin-susceptible population previously reported in SEA. These patterns confirm that the SEA population structure is driven substantially by drug resistance and positive selection.

Comparing pre- and post-peak-removal inferences, $N_e$ trajectories are tightly overlapping: pre- and post-removal point estimates sit within each other's $95\%$ confidence intervals across the full time range. Community assignments are also largely stable, with similar community sizes and adjusted Rand indices to the pre-correction solution; hierarchical clustering recovers the same C0/C2/C3/C4 vs.\ C1 grouping. In SEA, selection correction is therefore neither necessary nor beneficial for IBD-based inferences of $N_e$ and structure.

\subsection{IBD peak removal matters in WAF: high transmission, low background relatedness}

In WAF, peak removal produces materially different and significantly finer inferences, matching the pattern our simulations predict when background relatedness is low. This outcome is at odds with the null effect of peak removal in SEA, and we therefore tested whether the discrepancy is explained by baseline IBD. We hypothesised that the discrepancy stems from the high baseline genetic relatedness of SEA parasites, which dominates the IBD distribution and dilutes the contribution of localised selection peaks. We tested this hypothesis in two complementary ways—first by modifying the simulation to control background relatedness, and second by running the full empirical pipeline on WAF. First, we modified our simulations to condition the offspring generation process on parental pedigree relatedness for a subset of offspring candidates, thereby elevating background IBD sharing. In this high-relatedness regime, peak removal raised $N_e$ above the neutral expectation, suggesting over-correction, and failed to recover underlying population structure under the stepping-stone model. In the low-relatedness regime, peak removal closely tracked the neutral $N_e$ and recovered community labels as before.

Second, we analysed the WAF empirical dataset with matched unrelated-isolate criteria. In WAF, removing IBD peaks produced larger $N_e$ estimates in the most recent generations, with non-overlapping $95\%$ confidence intervals around $20$ generations ago. InfoMap community detection recovered markedly finer structure after peak removal: before correction, most isolates collapsed into a single dominant community; after correction, these isolates were redistributed across many smaller communities, mirroring the pattern from the stepping-stone simulation. Jackknife resampling confirmed that the change in community assignments is statistically significant. Figure~\ref{fig:correction_effect} summarises the transmission-setting contrast.

\begin{figure}[t]
  \centering
  \includegraphics[width=0.85\linewidth]{figures/fig4_correction_effect.pdf}
  \caption{Transmission setting modulates the effect of IBD peak removal on $N_e$ inference. In SEA, pre- and post-removal point estimates sit within each other's $95\%$ confidence intervals. In WAF, peak removal produces larger recent-generation $N_e$ estimates with non-overlapping CIs around $20$ generations ago.}
  \label{fig:correction_effect}
\end{figure}

Taken together, these results demonstrate that background genetic relatedness is a key modifier of selection-induced bias in IBD-based inference. Populations with low background relatedness, such as WAF, are susceptible to selection bias and benefit from peak-region removal. Populations with high background relatedness, such as SEA, are relatively insensitive to selection bias and risk over-correction if peak removal is applied.

\section{Discussion}

\textbf{Main finding.} \textit{P. falciparum} parasites have undergone strong, multi-locus selective sweeps driven by resistance to artemisinin, chloroquine, piperaquine, and antifolates. Our simulations and field-isolate analyses demonstrate that such sweeps distort the IBD distribution, depress recent $N_e$ estimates, and blur IBD-based population structure. A peak-removal correction, which excises IBD segments falling within validated sweep regions, mitigates this bias. Importantly, the benefit of correction is setting-dependent: it is substantial in high-transmission, low-relatedness populations such as those in sub-Saharan Africa, but minimal and potentially harmful in low-transmission, high-relatedness populations such as those in Southeast Asia.

\textbf{Mechanism.} Selection elongates IBD segments around selected loci, which increases both the long-segment tail of the segment-length distribution and pairwise total IBD. Because estimators such as \texttt{IBDNe} rely on a direct mapping from segment length to TMRCA, elongation biases $N_e$ downward in recent generations. This mechanism is orthogonal to changes in the background coalescent rate and is therefore invisible to methods that treat IBD as a neutral summary. At the network level, elongated selected-region segments inflate pairwise total IBD between distantly related pairs in unrelated demes, raising inter-population to intra-population IBD ratios and collapsing InfoMap communities. Our peak-removal correction attacks both effects simultaneously by deleting the segments most likely to carry the selection signature.

\textbf{Background relatedness as a modifier.} A striking finding of our study is that the benefit of correction depends on a scalar summary of the dataset: the genome-wide baseline IBD level. When baseline IBD is low, each selected-region segment contributes a large fractional increment to pairwise total IBD, and removing these segments reveals the underlying demographic signal. When baseline IBD is high, as in SEA after decades of declining transmission, pairwise total IBD is dominated by genome-wide sharing from a small effective population rather than by localised sweep-driven excess. Removing peak segments then shaves off a relatively small contribution, and can in fact over-correct by cutting long segments into shorter ones, shifting the distribution in ways that spuriously inflate $N_e$. This cautions against the automatic application of correction to every \textit{P. falciparum} dataset.

\textbf{Comparison with LD-based results.} Our conclusions appear to contradict those of \citet{novo2022ld}, who reported that LD-based $N_e$ estimators are essentially robust to natural selection. The apparent contradiction is, we believe, an artefact of simulation regime rather than a substantive disagreement. Their study used small selection coefficients, random placement of selected loci, and long evolutionary time scales—conditions under which the aggregate LD decay curve is only weakly perturbed. Our study models recent, strong, locus-anchored selection with establishment conditioning, which is the regime relevant for \textit{P. falciparum} under drug pressure. We also note that LD- and IBD-based estimators exploit different features of the haplotype distribution: LD integrates over pairs of sites at prescribed distances, while IBD depends on the tail of the segment-length distribution. Even if LD-based $N_e$ is robust to selection in a given regime, IBD-based $N_e$ may not be, and the two should not be conflated.

\textbf{Implications for genomic surveillance.} Practically, our results imply a simple decision rule: when IBD-based inferences are applied to high-transmission, low-relatedness populations, the analyst should remove IBD segments falling in validated peak regions before estimating $N_e$ or community structure. Peaks can be identified using threshold methods on the IBD coverage profile and validated with IBD-based selection statistics such as $X_{iR,s}$; in empirical data, such peaks co-localise with known drug-resistance genes including \textit{pfcrt}, \textit{pfmdr1}, \textit{dhps}, \textit{kelch13}, and \textit{arps10}, providing an independent check. In low-transmission, high-relatedness settings, correction should be applied cautiously or omitted, and the analyst should test for over-correction by examining whether $N_e$ estimates exceed reasonable neutral expectations. These recommendations are directly actionable for ongoing genomic surveillance programmes.

\textbf{A benchmarking framework for IBD callers.} Our simulation and true-IBD pipeline provides a foundation for evaluating IBD-calling tools under realistic \textit{P. falciparum} parameters, including strong selection, high recombination, and low effective SNP density. Prior validation of \texttt{hmmIBD} and \texttt{isoRelate} has relied primarily on parent--offspring or pedigree-based simulations, which do not reproduce the interaction of multi-locus selection with low effective SNP density. Our framework, built on \texttt{msprime}, \texttt{SLiM}, \texttt{tskit}, and the true-IBD inference algorithm \texttt{tskibd}, naturally extends to non-human, high-recombination species and to head-to-head benchmarking of new callers against ground truth.

\textbf{Limitations.} Several caveats accompany our analyses. First, \textit{P. falciparum} evolutionary parameters are not all well characterised, so we made assumptions about recombination rate, mutation rate, and generation time; biased assumptions would affect the quantitative level of simulated bias though not its qualitative direction. Second, our selection simulations place a sweep on every chromosome simultaneously to approximate multi-locus resistance, which may overestimate bias relative to the sequential sweeps seen empirically as drug policies change. Third, the empirical datasets are heterogeneous in sampling location and time, which can complicate temporal interpretation of $N_e$ and violate the panmixia assumption of \texttt{IBDNe}; restricting analyses to smaller time windows or finer geographic units is a natural remedy but trades off variance. Fourth, our peak-removal procedure uses a simple threshold with a boundary-extension heuristic; more sophisticated approaches that couple peak detection with explicit models of linked selection could yield tighter corrections and clearer criteria for when over-correction becomes a risk.

\textbf{Conclusion.} We have shown that strong recent positive selection biases IBD-based inferences of effective population size and population structure in \textit{P. falciparum}, and that a simple peak-region removal corrects this bias in low-background-relatedness settings but not in high-relatedness ones. Background relatedness is therefore a key decision variable for genomic surveillance. Our true-IBD simulation framework supports both this diagnostic and a broader programme of IBD-caller benchmarking for species with extreme evolutionary parameters. As \textit{P. falciparum} continues to evolve under sustained drug pressure, and as genomic surveillance expands geographically, transparent treatment of selection bias will be essential to accurate demographic and epidemiologic inference.

\bibliographystyle{plainnat}
\bibliography{references}

\end{document}
